\documentclass{beamer}
\usepackage[utf8]{inputenc}

\usetheme{Madrid}
\usecolortheme{default}
\usepackage{amsmath,amssymb,amsfonts,amsthm}
\usepackage{txfonts}
\usepackage{tkz-euclide}
\usepackage{listings}
\usepackage{adjustbox}
\usepackage{array}
\usepackage{tabularx}
\usepackage{gvv}
\usepackage{lmodern}
\usepackage{circuitikz}
\usepackage{tikz}
\usepackage{graphicx}

\setbeamertemplate{page number in head/foot}[totalframenumber]

\usepackage{tcolorbox}
\tcbuselibrary{minted,breakable,xparse,skins}



\definecolor{bg}{gray}{0.95}
\DeclareTCBListing{mintedbox}{O{}m!O{}}{%
	breakable=true,
	listing engine=minted,
	listing only,
	minted language=#2,
	minted style=default,
	minted options={%
		linenos,
		gobble=0,
		breaklines=true,
		breakafter=,,
		fontsize=\small,
		numbersep=8pt,
		#1},
	boxsep=0pt,
	left skip=0pt,
	right skip=0pt,
	left=25pt,
	right=0pt,
	top=3pt,
	bottom=3pt,
	arc=5pt,
	leftrule=0pt,
	rightrule=0pt,
	bottomrule=2pt,
	toprule=2pt,
	colback=bg,
	colframe=orange!70,
	enhanced,
	overlay={%
		\begin{tcbclipinterior}
			\fill[orange!20!white] (frame.south west) rectangle ([xshift=20pt]frame.north west);
	\end{tcbclipinterior}},
	#3,
}
\lstset{
	language=C,
	basicstyle=\ttfamily\small,
	keywordstyle=\color{blue},
	stringstyle=\color{orange},
	commentstyle=\color{green!60!black},
	numbers=left,
	numberstyle=\tiny\color{gray},
	breaklines=true,
	showstringspaces=false,
}
%------------------------------------------------------------
%This block of code defines the information to appear in the
%Title page
\title %optional
{12.580}
%\subtitle{A short story}

\author % (optional)
{RAVULA SHASHANK REDDY - EE25BTECH11047}

 \begin{document}
	
	
	\frame{\titlepage}
	\begin{frame}{Question}
    Let \(\vec{V}\) be the vector space of all \(3 \times 3\) matrices with complex entries over the real field. If  
\[
\vec{W_1} = \{\vec{A} \in \vec{V} : \vec{A} = \overline{\vec{A}}^\top \} \quad \text{and} \quad \vec{W_2} = \{\vec{A} \in \vec{V} : \mathrm{trace}(\vec{A}) = 0\},
\]  
then find \(\dim(\vec{W_1} + \vec{W_2})\).
\end{frame}
\begin{frame}{Solution}
    \begin{align}
\vec{V} &= \text{all } 3 \times 3 \text{ complex matrices over } \mathbb{R} 
\end{align}

For one complex element 2 real parameters are required.
\begin{align}
\dim_\mathbb{R}(\vec{V}) &= 3 \times 3 \times 2 = 18
\end{align}

\begin{align}
\vec{W_1} &= \{\vec{A} \in \vec{V}: \vec{A} = \overline{\vec{A}}^\top \} \\
\vec{A} = \myvec{a_{11} & a_{12} & a_{13} \\
\overline{a_{12}}& a_{22} & a_{23} \\
\overline{a_{13}} &\overline{a_{23}}& a_{33} }
a_{ii} &\in \mathbb{R}, \quad a_{ij} \in \mathbb{C}, \; i<j \\
\dim_\mathbb{R}(\vec{W_1}) &= 3 + 3 \times 2 = 9
\end{align}

\begin{align}
\vec{W_2} &= \{\vec{A} \in \vec{V}: \mathrm{trace}(\vec{A}) = 0\} 
\end{align}
\end{frame}
\begin{frame}{Solution}
So one element along diagonal has constraint.
\begin{align}
\dim_\mathbb{R}(\vec{W_2}) &= ((3 \times 3)-1) \times 2  = 16
\end{align}

\begin{align}
\dim_\mathbb{R}(\vec{W_1} \cap \vec{W_2}) &= 9 - 1 = 8
\end{align}

\begin{align}
\dim(\vec{W_1} + \vec{W_2}) &= \dim(\vec{W_1}) + \dim(\vec{W_2}) - \dim(\vec{W_1} \cap \vec{W_2}) \\
&= 9 + 16 - 8 \\
&= 17
\end{align}
\end{frame}
\begin{frame}[fragile]
\frametitle{C Code}
\begin{lstlisting}
    #include <stdio.h>

int main() {
    int n = 3; // 3x3 matrix
    int complex_entries = n * n;  // total complex entries
    
    // Step 1: Dimension of V (all complex entries over R)
    int dimV = complex_entries * 2; // 2 real numbers per complex entry

    // Step 2: Dimension of W1 (Hermitian)
    // Diagonal entries: n real numbers
    // Off-diagonal entries: n*(n-1)/2 complex numbers = (n*(n-1)) real
    int dimW1 = n + (n*(n-1)); // 3 + 6 = 9
\end{lstlisting}
\end{frame}
\begin{frame}[fragile]
\frametitle{C Code}
\begin{lstlisting}
    // Step 3: Dimension of W2 (trace zero)
    // Trace=0 gives 1 complex constraint = 2 real
    int dimW2 = dimV - 2; // 18 - 2 = 16

    // Step 4: Dimension of W1 ∩ W2
    int dimIntersection = dimW1 - 1; // 1 real reduction from trace constraint

    // Step 5: Dimension of W1 + W2
    int dimSum = dimW1 + dimW2 - dimIntersection;

    printf("%d\n", dimSum);

    return 0;
}
\end{lstlisting}
\end{frame}
\begin{frame}[fragile]
\frametitle{Python Direct}
\begin{lstlisting}
    import numpy as np

def dimension_W1_plus_W2(n=3):
    # Dimension of V: all n x n complex matrices over R
    dimV = n * n * 2  # 2 real numbers per complex entry

    # Dimension of W1 (Hermitian)
    # Diagonal: n real numbers
    # Off-diagonal: n*(n-1)/2 complex numbers = n*(n-1) real numbers
    dimW1 = n + n*(n-1)

    # Dimension of W2 (trace zero)
    # Trace=0 gives 1 complex constraint = 2 real
    dimW2 = dimV - 2
\end{lstlisting}
\end{frame}
\begin{frame}[fragile]
\frametitle{Python Direct}
\begin{lstlisting}
    # Dimension of W1 ∩ W2
    dimIntersection = dimW1 - 1  # trace constraint reduces 1 real dimension

    # Dimension of W1 + W2
    dimSum = dimW1 + dimW2 - dimIntersection
    return dimSum

if __name__ == "__main__":
    n = 3  # 3x3 matrix
    print(dimension_W1_plus_W2(n))
\end{lstlisting}
\end{frame}

\end{document}
